\section{Introduction}

The fourth cohomological moment in the H\'enon source program carries an
apparently awkward half-power over the quadratic field
\(K=\Q(\rho)\).  A direct \(K\)-compatible system cannot realize that
half-power without acquiring a nonintegral rank.  The obstruction leaves
one algebraic escape: the two \(K\)-places above a split rational prime
could be the conjugate copies of a single object over \(\Q\).

This paper constructs that rational object.  The construction begins at the
equation level, before any local factor is manipulated.  Reversal of the
source chronology is accompanied by alternating powers of \(\rho\); the
resulting monomial map is a Weil descent datum.  Solving its fixed-vector
equations gives one rational coordinate system in every row and preserves
the weighted closing edge.  Thus the repair is geometric rather than a
formal choice of a square-root branch.

The fourth row contains more structure.  The preceding source analysis
found an order-\(24\) group
\(\Dih(C_{12})\) and a rank-\(10\) Reynolds summand in the middle
cohomology.  Conjugation by the rational descent does not fix the
twenty-four transformations individually.  Instead it acts by
\[
r\longmapsto r^{-1},\qquad s\longmapsto sr^{-1}.
\]
The correct rational object is consequently a nonconstant finite \'etale
group scheme split by \(K\).  Its graph sum is invariant, so rational
restriction and corestriction descend the Chow projector.

\subsection*{Contributions}

The paper proves four claims.

\begin{enumerate}[leftmargin=*]
\item For every \(n\ge2\), we construct an explicit \(\Q\)-model of the
source-ordered \((2,3)\) intersection and prove a closed determinant
formula for the coordinate change.  This statement concerns equations;
smoothness and motives are not promoted beyond the certified rows.
\item For \(n=2,3,4\), we descend the full cohomological packet of rank
\(4^n-1\).  At good split primes, its exponent changes from \(2/n\) over
\(K\) to \(4/n\) over \(\Q\).  The \(n=4\) half-power clears exactly.
\item We descend the order-\(24\) Reynolds correspondence and obtain
rational Chow summands of ranks \(10\) and \(158\).  The rank-\(10\)
summand has Hodge ledger \((1,4,4,1)\) and becomes of
Calabi--Yau-threefold Hodge type after one Tate twist.
\item At good primes, the raw rank-\(10\) Frobenius polynomial is
\(\ell\)-independent, belongs to \(\Z[T]\), has weight \(5\), and obeys
the expected degree-five reciprocity.  This strict compatibility statement
does not use semisimplicity.
\end{enumerate}

\subsection*{Relation to existing geometry}

Favero--Iliev--Katzarkov give the Hodge numbers and Cayley-ring description
for smooth \((2,3)\) fivefolds in \(\PP^7\)
\cite[\S5.1, equation (6), and \S5.4, equation (8)]{FIK2012}.
Those results supply context for the ledger, not the source-specific
descent or projector.  Laterveer records ordinary nonmiddle
Chow--K\"unneth projectors for intersections of a quadric and a cubic
\cite[proof of Theorem 4.1]{Laterveer2021}; the multiplicative theorem
there assumes even ambient dimension and does not apply to \(\PP^7\).
Vial gives broader finite-dimensional-motive context
\cite[Example 4.12]{Vial2013}.  None of these references identifies the
H\'enon-derived rational form or its twisted dihedral core.

Our descent proof is explicit and therefore avoids a general effectivity
claim for arbitrary projective descent data.  For Chow cycles, we use the
standard proper-pushforward and flat-pullback degree formulas
\cite[\S\S1.4, 1.7]{Fulton1998}.  Kahn's motivic descent theorem
\cite[Theorems 1 and 7.1]{Kahn2023} gives a modern categorical backdrop,
but the graph-sum argument below is elementary with rational coefficients.
For local polynomials, the comparison input is the
correspondence-trace theorem of Katz--Messing
\cite[Theorem 2(2)]{KatzMessing1974}.

\subsection*{Scope}

Every Frobenius in this paper is geometric, normalized by
\(\Frob_p\mid\Q_\ell(-1)=p\).  The rational exponent clearing is local to
good split primes.  At inert primes, Frobenius eigenvalues are squared
rather than duplicated, which blocks a global half-root.  We do not claim
smoothness for \(n\ge5\), automorphy, meromorphic continuation, a functional
equation, a Calabi--Yau realization, or Chow indecomposability.  The exact
irreducibility gate for the rank-\(10\) polynomial is recorded only as a
successor problem.
